% --- Archon physics formalization source begin ---
% archon:physics
% archon:covers PhyXMiniProblems/problem_phyx_mini_0986.lean
% archon:source-report reports/phyx_mini/problem_phyx_mini_0986.source.json

\chapter{Physics problem phyx\_mini\_0986}
\label{ch:PhyXMiniProblems_problem_phyx_mini_0986}

\paragraph{Problem source.}
\#\# Physical scenario

In the circuit, the switch \textbackslash{}( S \textbackslash{}) has been open for a long time and is suddenly closed. Neither the battery nor the inductors have any appreciable resistance.

\#\# Question

What does the voltmeter read 0.115 ms after \textbackslash{}( S \textbackslash{}) is closed?

\#\# Answer choices

- A: A: 0V

- B: B: 9.0V

- C: C: 20.0V

- D: D: 11.0V

\#\# Figure caption (auxiliary; use the image as primary evidence)

The image depicts an electrical circuit diagram containing the following components:

1. **Voltage Source**: A 20.0 V battery is shown on the left side, indicating the source of electrical power.

2. **Ammeter (A)**: Positioned at the bottom left, it measures the current flowing through the circuit.

3. **Switch (S)**: Located near the top left, it controls the flow of current in the circuit.

4. **Resistors**:
   - A 50.0 Ω resistor is connected in series with the switch.
   - A 25.0 Ω resistor is connected in series with the ammeter.

5. **Voltmeter (V)**: Positioned in the center, it measures the voltage across certain components.

6. **Inductors**:
   - A 12.0 mH inductor is connected near the top right in parallel with the voltmeter.
   - An 18.0 mH inductor is connected above another inductor.
   - A 15.0 mH inductor is positioned toward the right, connected in series with the 18.0 mH inductor.

The circuit is organized with a combination of series and parallel components, demonstrating various electrical relationships and configurations.

\#\# Dataset metadata

Electric Circuits; Physical Model Grounding Reasoning, Numerical Reasoning

\paragraph{Recorded answer/context.}
B: B: 9.0V

\paragraph{Figure/image path.}
/root/proposal\_for\_physic/science-mango/phyx\_mini\_run/phyx\_data/test\_image/986.png

\paragraph{Formalization target.}
create a compiling Lean file with sorry bodies at `PhyXMiniProblems/problem\_phyx\_mini\_0986.lean`. The Lean declarations must preserve the physical quantities, dimensions or dimensional roles, figure labels, governing-law hypotheses, and final relation expressed by this problem.
Use Mathlib/Physlib names found through LeanExplore where available. If a physics API is missing, introduce faithful local abstractions rather than scalar placeholder aliases.

\begin{theorem}[Physics formalization target]
\label{thm:physics:phyx_mini_0986:target}
\lean{PhyXMiniProblems.ProblemPhyXMini0986.voltmeter_reads_nine_volts_at_point115_milliseconds}
\uses{def:physics:phyx-mini-0986:phyxminiproblems-problemphyxmini0986-switchedinductorcircuit, def:physics:phyx-mini-0986:phyxminiproblems-problemphyxmini0986-matchesprimarycircuitfigure, def:physics:phyx-mini-0986:phyxminiproblems-problemphyxmini0986-matchescircuitnumericaldata, def:physics:phyx-mini-0986:phyxminiproblems-problemphyxmini0986-hasidealcomponentphysics, def:physics:phyx-mini-0986:phyxminiproblems-problemphyxmini0986-isopenlongthenclosed, def:physics:phyx-mini-0986:phyxminiproblems-problemphyxmini0986-satisfiesidealrltransientlaws, def:physics:phyx-mini-0986:phyxminiproblems-problemphyxmini0986-hasonedecimalvoltmeterdisplay}
The assigned autoformalize agent should translate this physics problem into Lean declarations in the covered file, with theorem and lemma proof bodies written as `by sorry`.
\end{theorem}
\begin{proof}
This is an autoformalization task, not a proof task. Produce statements that can later be proved by the physics prover without weakening the physical model.
\end{proof}

% --- Archon declaration topology begin ---
\section{Declaration topology}
\begin{definition}[voltage Dimension]
  \label{def:physics:phyx-mini-0986:phyxminiproblems-problemphyxmini0986-voltagedimension}
  \lean{PhyXMiniProblems.ProblemPhyXMini0986.voltageDimension}
  The dimension M L² T⁻² C⁻¹ of electric potential difference
\end{definition}
\begin{proof}
  This is the declared model definition of voltage Dimension.
\end{proof}
\begin{definition}[electric Current Dimension]
  \label{def:physics:phyx-mini-0986:phyxminiproblems-problemphyxmini0986-electriccurrentdimension}
  \lean{PhyXMiniProblems.ProblemPhyXMini0986.electricCurrentDimension}
  The dimension C T⁻¹ of electric current
\end{definition}
\begin{proof}
  This is the declared model definition of electric Current Dimension.
\end{proof}
\begin{definition}[electrical Resistance Dimension]
  \label{def:physics:phyx-mini-0986:phyxminiproblems-problemphyxmini0986-electricalresistancedimension}
  \lean{PhyXMiniProblems.ProblemPhyXMini0986.electricalResistanceDimension}
  The dimension M L² T⁻¹ C⁻² of electrical resistance
\end{definition}
\begin{proof}
  This is the declared model definition of electrical Resistance Dimension.
\end{proof}
\begin{definition}[inductance Dimension]
  \label{def:physics:phyx-mini-0986:phyxminiproblems-problemphyxmini0986-inductancedimension}
  \lean{PhyXMiniProblems.ProblemPhyXMini0986.inductanceDimension}
  The dimension M L² C⁻² of self-inductance
\end{definition}
\begin{proof}
  This is the declared model definition of inductance Dimension.
\end{proof}
\begin{definition}[Voltage Quantity]
  \label{def:physics:phyx-mini-0986:phyxminiproblems-problemphyxmini0986-voltagequantity}
  \lean{PhyXMiniProblems.ProblemPhyXMini0986.VoltageQuantity}
  \uses{def:physics:phyx-mini-0986:phyxminiproblems-problemphyxmini0986-voltagedimension}
  A signed, unit-independent voltage
\end{definition}
\begin{proof}
  This is the declared model definition of Voltage Quantity.
\end{proof}
\begin{definition}[Electric Current Quantity]
  \label{def:physics:phyx-mini-0986:phyxminiproblems-problemphyxmini0986-electriccurrentquantity}
  \lean{PhyXMiniProblems.ProblemPhyXMini0986.ElectricCurrentQuantity}
  \uses{def:physics:phyx-mini-0986:phyxminiproblems-problemphyxmini0986-electriccurrentdimension}
  A signed, unit-independent electric current
\end{definition}
\begin{proof}
  This is the declared model definition of Electric Current Quantity.
\end{proof}
\begin{definition}[Resistance Magnitude]
  \label{def:physics:phyx-mini-0986:phyxminiproblems-problemphyxmini0986-resistancemagnitude}
  \lean{PhyXMiniProblems.ProblemPhyXMini0986.ResistanceMagnitude}
  \uses{def:physics:phyx-mini-0986:phyxminiproblems-problemphyxmini0986-electricalresistancedimension}
  A nonnegative, unit-independent electrical resistance
\end{definition}
\begin{proof}
  This is the declared model definition of Resistance Magnitude.
\end{proof}
\begin{definition}[Inductance Magnitude]
  \label{def:physics:phyx-mini-0986:phyxminiproblems-problemphyxmini0986-inductancemagnitude}
  \lean{PhyXMiniProblems.ProblemPhyXMini0986.InductanceMagnitude}
  \uses{def:physics:phyx-mini-0986:phyxminiproblems-problemphyxmini0986-inductancedimension}
  A nonnegative, unit-independent self-inductance
\end{definition}
\begin{proof}
  This is the declared model definition of Inductance Magnitude.
\end{proof}
\begin{definition}[Time Quantity]
  \label{def:physics:phyx-mini-0986:phyxminiproblems-problemphyxmini0986-timequantity}
  \lean{PhyXMiniProblems.ProblemPhyXMini0986.TimeQuantity}
  A signed, unit-independent time coordinate relative to switch closing
\end{definition}
\begin{proof}
  This is the declared model definition of Time Quantity.
\end{proof}
\begin{definition}[voltage In Volts]
  \label{def:physics:phyx-mini-0986:phyxminiproblems-problemphyxmini0986-voltageinvolts}
  \lean{PhyXMiniProblems.ProblemPhyXMini0986.voltageInVolts}
  \uses{def:physics:phyx-mini-0986:phyxminiproblems-problemphyxmini0986-voltagequantity}
  Coherent-SI voltage readout, in volts
\end{definition}
\begin{proof}
  This is the declared model definition of voltage In Volts.
\end{proof}
\begin{definition}[current In Amperes]
  \label{def:physics:phyx-mini-0986:phyxminiproblems-problemphyxmini0986-currentinamperes}
  \lean{PhyXMiniProblems.ProblemPhyXMini0986.currentInAmperes}
  \uses{def:physics:phyx-mini-0986:phyxminiproblems-problemphyxmini0986-electriccurrentquantity}
  Coherent-SI current readout, in amperes
\end{definition}
\begin{proof}
  This is the declared model definition of current In Amperes.
\end{proof}
\begin{definition}[resistance In Ohms]
  \label{def:physics:phyx-mini-0986:phyxminiproblems-problemphyxmini0986-resistanceinohms}
  \lean{PhyXMiniProblems.ProblemPhyXMini0986.resistanceInOhms}
  \uses{def:physics:phyx-mini-0986:phyxminiproblems-problemphyxmini0986-resistancemagnitude}
  Coherent-SI resistance readout, in ohms
\end{definition}
\begin{proof}
  This is the declared model definition of resistance In Ohms.
\end{proof}
\begin{definition}[inductance In Henries]
  \label{def:physics:phyx-mini-0986:phyxminiproblems-problemphyxmini0986-inductanceinhenries}
  \lean{PhyXMiniProblems.ProblemPhyXMini0986.inductanceInHenries}
  \uses{def:physics:phyx-mini-0986:phyxminiproblems-problemphyxmini0986-inductancemagnitude}
  Coherent-SI inductance readout, in henries
\end{definition}
\begin{proof}
  This is the declared model definition of inductance In Henries.
\end{proof}
\begin{definition}[inductance In Millihenries]
  \label{def:physics:phyx-mini-0986:phyxminiproblems-problemphyxmini0986-inductanceinmillihenries}
  \lean{PhyXMiniProblems.ProblemPhyXMini0986.inductanceInMillihenries}
  \uses{def:physics:phyx-mini-0986:phyxminiproblems-problemphyxmini0986-inductancemagnitude, def:physics:phyx-mini-0986:phyxminiproblems-problemphyxmini0986-inductanceinhenries}
  Inductance readout in millihenries
\end{definition}
\begin{proof}
  This is the declared model definition of inductance In Millihenries.
\end{proof}
\begin{definition}[time In Seconds]
  \label{def:physics:phyx-mini-0986:phyxminiproblems-problemphyxmini0986-timeinseconds}
  \lean{PhyXMiniProblems.ProblemPhyXMini0986.timeInSeconds}
  \uses{def:physics:phyx-mini-0986:phyxminiproblems-problemphyxmini0986-timequantity}
  Coherent-SI time readout, in seconds
\end{definition}
\begin{proof}
  This is the declared model definition of time In Seconds.
\end{proof}
\begin{definition}[time In Milliseconds]
  \label{def:physics:phyx-mini-0986:phyxminiproblems-problemphyxmini0986-timeinmilliseconds}
  \lean{PhyXMiniProblems.ProblemPhyXMini0986.timeInMilliseconds}
  \uses{def:physics:phyx-mini-0986:phyxminiproblems-problemphyxmini0986-timequantity, def:physics:phyx-mini-0986:phyxminiproblems-problemphyxmini0986-timeinseconds}
  Time readout in milliseconds
\end{definition}
\begin{proof}
  This is the declared model definition of time In Milliseconds.
\end{proof}
\begin{definition}[Circuit Node]
  \label{def:physics:phyx-mini-0986:phyxminiproblems-problemphyxmini0986-circuitnode}
  \lean{PhyXMiniProblems.ProblemPhyXMini0986.CircuitNode}
  Electrically distinct nodes visible in image 986.png
\end{definition}
\begin{proof}
  This is the declared model definition of Circuit Node.
\end{proof}
\begin{definition}[Component Label]
  \label{def:physics:phyx-mini-0986:phyxminiproblems-problemphyxmini0986-componentlabel}
  \lean{PhyXMiniProblems.ProblemPhyXMini0986.ComponentLabel}
  Every labeled circuit element visible in image 986.png
\end{definition}
\begin{proof}
  This is the declared model definition of Component Label.
\end{proof}
\begin{definition}[Circuit Figure]
  \label{def:physics:phyx-mini-0986:phyxminiproblems-problemphyxmini0986-circuitfigure}
  \lean{PhyXMiniProblems.ProblemPhyXMini0986.CircuitFigure}
  \uses{def:physics:phyx-mini-0986:phyxminiproblems-problemphyxmini0986-circuitnode, def:physics:phyx-mini-0986:phyxminiproblems-problemphyxmini0986-componentlabel}
  The endpoint data read directly from the supplied circuit diagram
\end{definition}
\begin{proof}
  This is the declared model definition of Circuit Figure.
\end{proof}
\begin{definition}[Switched Inductor Circuit]
  \label{def:physics:phyx-mini-0986:phyxminiproblems-problemphyxmini0986-switchedinductorcircuit}
  \lean{PhyXMiniProblems.ProblemPhyXMini0986.SwitchedInductorCircuit}
  \uses{def:physics:phyx-mini-0986:phyxminiproblems-problemphyxmini0986-voltagequantity, def:physics:phyx-mini-0986:phyxminiproblems-problemphyxmini0986-electriccurrentquantity, def:physics:phyx-mini-0986:phyxminiproblems-problemphyxmini0986-resistancemagnitude, def:physics:phyx-mini-0986:phyxminiproblems-problemphyxmini0986-inductancemagnitude, def:physics:phyx-mini-0986:phyxminiproblems-problemphyxmini0986-timequantity, def:physics:phyx-mini-0986:phyxminiproblems-problemphyxmini0986-circuitfigure}
  The independent physical quantities and time histories of the experiment. The equivalent resistance and inductances are included as figure-derived physical quantities; their reduction equations are laws rather than definitions
\end{definition}
\begin{proof}
  This is the declared model definition of Switched Inductor Circuit.
\end{proof}
\begin{definition}[Matches Primary Circuit Figure]
  \label{def:physics:phyx-mini-0986:phyxminiproblems-problemphyxmini0986-matchesprimarycircuitfigure}
  \lean{PhyXMiniProblems.ProblemPhyXMini0986.MatchesPrimaryCircuitFigure}
  \uses{def:physics:phyx-mini-0986:phyxminiproblems-problemphyxmini0986-switchedinductorcircuit}
  Exact component incidences transcribed from the primary raster
\end{definition}
\begin{proof}
  This is the declared model definition of Matches Primary Circuit Figure.
\end{proof}
\begin{definition}[Matches Circuit Numerical Data]
  \label{def:physics:phyx-mini-0986:phyxminiproblems-problemphyxmini0986-matchescircuitnumericaldata}
  \lean{PhyXMiniProblems.ProblemPhyXMini0986.MatchesCircuitNumericalData}
  \uses{def:physics:phyx-mini-0986:phyxminiproblems-problemphyxmini0986-voltageinvolts, def:physics:phyx-mini-0986:phyxminiproblems-problemphyxmini0986-resistanceinohms, def:physics:phyx-mini-0986:phyxminiproblems-problemphyxmini0986-inductanceinmillihenries, def:physics:phyx-mini-0986:phyxminiproblems-problemphyxmini0986-timeinseconds, def:physics:phyx-mini-0986:phyxminiproblems-problemphyxmini0986-timeinmilliseconds, def:physics:phyx-mini-0986:phyxminiproblems-problemphyxmini0986-switchedinductorcircuit}
  All numerical labels shown in the problem and its primary figure
\end{definition}
\begin{proof}
  This is the declared model definition of Matches Circuit Numerical Data.
\end{proof}
\begin{definition}[Has Ideal Component Physics]
  \label{def:physics:phyx-mini-0986:phyxminiproblems-problemphyxmini0986-hasidealcomponentphysics}
  \lean{PhyXMiniProblems.ProblemPhyXMini0986.HasIdealComponentPhysics}
  \uses{def:physics:phyx-mini-0986:phyxminiproblems-problemphyxmini0986-currentinamperes, def:physics:phyx-mini-0986:phyxminiproblems-problemphyxmini0986-resistanceinohms, def:physics:phyx-mini-0986:phyxminiproblems-problemphyxmini0986-switchedinductorcircuit}
  The negligible-resistance assumptions stated in the prose, together with the standard ideal-meter behavior used by the circuit diagram
\end{definition}
\begin{proof}
  This is the declared model definition of Has Ideal Component Physics.
\end{proof}
\begin{definition}[Is Open Long Then Closed]
  \label{def:physics:phyx-mini-0986:phyxminiproblems-problemphyxmini0986-isopenlongthenclosed}
  \lean{PhyXMiniProblems.ProblemPhyXMini0986.IsOpenLongThenClosed}
  \uses{def:physics:phyx-mini-0986:phyxminiproblems-problemphyxmini0986-currentinamperes, def:physics:phyx-mini-0986:phyxminiproblems-problemphyxmini0986-timeinseconds, def:physics:phyx-mini-0986:phyxminiproblems-problemphyxmini0986-switchedinductorcircuit}
  The switch history and zero-current state produced by being open a long time
\end{definition}
\begin{proof}
  This is the declared model definition of Is Open Long Then Closed.
\end{proof}
\begin{definition}[Satisfies Ideal RLTransient Laws]
  \label{def:physics:phyx-mini-0986:phyxminiproblems-problemphyxmini0986-satisfiesidealrltransientlaws}
  \lean{PhyXMiniProblems.ProblemPhyXMini0986.SatisfiesIdealRLTransientLaws}
  \uses{def:physics:phyx-mini-0986:phyxminiproblems-problemphyxmini0986-voltageinvolts, def:physics:phyx-mini-0986:phyxminiproblems-problemphyxmini0986-currentinamperes, def:physics:phyx-mini-0986:phyxminiproblems-problemphyxmini0986-resistanceinohms, def:physics:phyx-mini-0986:phyxminiproblems-problemphyxmini0986-inductanceinhenries, def:physics:phyx-mini-0986:phyxminiproblems-problemphyxmini0986-timeinseconds, def:physics:phyx-mini-0986:phyxminiproblems-problemphyxmini0986-switchedinductorcircuit}
  Kirchhoff laws, series/parallel reduction, and the standard closed-form voltage response for an ideal RL network after a DC step. The response law is uniform in time and in all component values; it contains neither the requested time nor the answer choice
\end{definition}
\begin{proof}
  This is the declared model definition of Satisfies Ideal RLTransient Laws.
\end{proof}
\begin{definition}[Has One Decimal Voltmeter Display]
  \label{def:physics:phyx-mini-0986:phyxminiproblems-problemphyxmini0986-hasonedecimalvoltmeterdisplay}
  \lean{PhyXMiniProblems.ProblemPhyXMini0986.HasOneDecimalVoltmeterDisplay}
  \uses{def:physics:phyx-mini-0986:phyxminiproblems-problemphyxmini0986-voltageinvolts, def:physics:phyx-mini-0986:phyxminiproblems-problemphyxmini0986-switchedinductorcircuit}
  Generic calibration law for a voltmeter displaying volts to one decimal place
\end{definition}
\begin{proof}
  This is the declared model definition of Has One Decimal Voltmeter Display.
\end{proof}
\begin{lemma}[total series resistance is 75 ohms]
  \label{lem:physics:phyx-mini-0986:phyxminiproblems-problemphyxmini0986-total-series-resistance-is-75-ohms}
  \lean{PhyXMiniProblems.ProblemPhyXMini0986.total_series_resistance_is_75_ohms}
  \uses{def:physics:phyx-mini-0986:phyxminiproblems-problemphyxmini0986-resistanceinohms, def:physics:phyx-mini-0986:phyxminiproblems-problemphyxmini0986-switchedinductorcircuit, def:physics:phyx-mini-0986:phyxminiproblems-problemphyxmini0986-matchescircuitnumericaldata, def:physics:phyx-mini-0986:phyxminiproblems-problemphyxmini0986-satisfiesidealrltransientlaws}
  The two visible series resistors have equivalent resistance 75 Ω
\end{lemma}
\begin{proof}
  The conclusion follows from the governing relations and figure readouts named in the statement of total series resistance is 75 ohms.
\end{proof}
\begin{lemma}[parallel equivalent inductance is 10 8 millihenries]
  \label{lem:physics:phyx-mini-0986:phyxminiproblems-problemphyxmini0986-parallel-equivalent-inductance-is-10-8-millihenries}
  \lean{PhyXMiniProblems.ProblemPhyXMini0986.parallel_equivalent_inductance_is_10_8_millihenries}
  \uses{def:physics:phyx-mini-0986:phyxminiproblems-problemphyxmini0986-inductanceinmillihenries, def:physics:phyx-mini-0986:phyxminiproblems-problemphyxmini0986-switchedinductorcircuit, def:physics:phyx-mini-0986:phyxminiproblems-problemphyxmini0986-matchescircuitnumericaldata, def:physics:phyx-mini-0986:phyxminiproblems-problemphyxmini0986-satisfiesidealrltransientlaws}
  The parallel combination 18 mH ∥ (12 mH + 15 mH) is 10.8 mH
\end{lemma}
\begin{proof}
  The conclusion follows from the governing relations and figure readouts named in the statement of parallel equivalent inductance is 10 8 millihenries.
\end{proof}
\begin{lemma}[observation voltage is within one millivolt of nine]
  \label{lem:physics:phyx-mini-0986:phyxminiproblems-problemphyxmini0986-observation-voltage-is-within-one-millivolt-of-nine}
  \lean{PhyXMiniProblems.ProblemPhyXMini0986.observation_voltage_is_within_one_millivolt_of_nine}
  \uses{def:physics:phyx-mini-0986:phyxminiproblems-problemphyxmini0986-voltageinvolts, def:physics:phyx-mini-0986:phyxminiproblems-problemphyxmini0986-switchedinductorcircuit, def:physics:phyx-mini-0986:phyxminiproblems-problemphyxmini0986-matchescircuitnumericaldata, def:physics:phyx-mini-0986:phyxminiproblems-problemphyxmini0986-isopenlongthenclosed, def:physics:phyx-mini-0986:phyxminiproblems-problemphyxmini0986-satisfiesidealrltransientlaws}
  At 0.115 ms, the physical voltage is within 0.001 V of 9 V
\end{lemma}
\begin{proof}
  The conclusion follows from the governing relations and figure readouts named in the statement of observation voltage is within one millivolt of nine.
\end{proof}
% --- Archon declaration topology end ---
% --- Archon physics formalization source end ---
